\documentclass[11pt,oneside]{article}

\pdfminorversion=7
\usepackage{QuizStyle}

\hypersetup{
    pdftitle={20261112 Quiz for MATH102 Calculus II},
    pdfauthor={Jungwoo Kim},
    pdfsubject={Quiz for MATH102 Calculus II},
    pdfcreator={LaTeX}
}

\begin{document}

\quiztitle{20261112 Quiz}

\begin{problem}{6.4}
Show that the ball
\[
    B=\{\mathbf X\in\R^3\mid\|\mathbf X\|\leq1\}
\]
is a smoothly bounded set.
\end{problem}

\begin{solution}
The set $B$ is closed and bounded, and its boundary is the unit sphere
\[
    \partial B=\{(x,y,z)\in\R^3\mid x^2+y^2+z^2=1\}.
\]
Put
\[
    D=\{(u,v)\in\R^2\mid u^2+v^2\leq\tfrac23\}.
\]
This disk is smoothly bounded in $\R^2$, since its boundary circle is a finite union of $C^1$ coordinate graphs.  On $D$, the functions
\[
    (u,v)\longmapsto\pm\sqrt{1-u^2-v^2}
\]
are continuously differentiable: indeed, $1-u^2-v^2\geq1/3>0$ on $D$.

Using these functions, form the six graphs
\[
\begin{aligned}
    x&=\pm\sqrt{1-y^2-z^2}, &&(y,z)\in D,\\
    y&=\pm\sqrt{1-x^2-z^2}, &&(x,z)\in D,\\
    z&=\pm\sqrt{1-x^2-y^2}, &&(x,y)\in D.
\end{aligned}
\]
Each graph lies in $\partial B$. Conversely, if $(x,y,z)\in\partial B$, then
\[
    \max\{|x|,|y|,|z|\}\geq\frac1{\sqrt3},
\]
because $x^2+y^2+z^2=1$. Thus at least one coordinate can be written as one of the six graphs above, with the other two coordinates belonging to $D$. Hence these six $C^1$ graphs cover $\partial B$.

Therefore $B$ is smoothly bounded by Definition~6.4.
\end{solution}

\begin{problem}{6.18}
Write a Riemann sum for
\[
    \iint_{[0,1]\times[0,1]}x^2y^3\,\dd A
\]
using the upper right corner of each $\tfrac14$-square as its sample point.
\end{problem}

\begin{solution}
Divide $[0,1]\times[0,1]$ into the $16$ squares
\[
    Q_{ij}=\left[\frac{i-1}{4},\frac i4\right]
           \times\left[\frac{j-1}{4},\frac j4\right],
    \qquad 1\leq i,j\leq4.
\]
Each square has area $\Delta A=1/16$, and its upper right corner is
$(i/4,j/4)$. Thus the requested Riemann sum is
\[
    \boxed{\sum_{i=1}^{4}\sum_{j=1}^{4}
       \left(\frac i4\right)^2\left(\frac j4\right)^3\frac1{16}}.
\]
\end{solution}

\begin{problem}{6.20}
Let $D=[a,b]\times[c,d]$, and let $f$ and $g$ be continuous on
$[a,b]$ and $[c,d]$, respectively. Use Riemann sums to prove that
\[
    \iint_D f(x)g(y)\,\dd A
    =\left(\int_a^b f(x)\,\dd x\right)
     \left(\int_c^d g(y)\,\dd y\right).
\]
\end{problem}

\begin{solution}
Let $h=2^{-n}$, as in the Riemann-sum construction of Section~6.2, and set
\[
\begin{aligned}
    I_h&=\{i\in\Z\mid[ih,(i+1)h]\subseteq[a,b]\},\\
    J_h&=\{j\in\Z\mid[jh,(j+1)h]\subseteq[c,d]\}.
\end{aligned}
\]
The $h$-squares contained in $D$ are exactly the products of the intervals
indexed by $I_h$ and $J_h$. Choose $x_i\in[ih,(i+1)h]$ for $i\in I_h$ and
$y_j\in[jh,(j+1)h]$ for $j\in J_h$. The corresponding Riemann sum separates:
\[
\begin{aligned}
    S(fg,D,h)
    &=\sum_{i\in I_h}\sum_{j\in J_h}f(x_i)g(y_j)h^2\\
    &=\left(\sum_{i\in I_h}f(x_i)h\right)
      \left(\sum_{j\in J_h}g(y_j)h\right).
\end{aligned}
\]
The intervals omitted at the ends of $[a,b]$ and $[c,d]$ have total lengths
at most $2h$ in each case. Since $f$ and $g$ are continuous and therefore
bounded, these omissions do not affect the limits of the one-variable
Riemann sums. Consequently,
\[
    \lim_{h\to0}\sum_{i\in I_h}f(x_i)h=\int_a^b f(x)\,\dd x,
    \qquad
    \lim_{h\to0}\sum_{j\in J_h}g(y_j)h=\int_c^d g(y)\,\dd y.
\]
Also, $f(x)g(y)$ is continuous on the rectangle $D$, so by the Riemann-sum
convergence result of Section~6.2 (including its extension to functions of
arbitrary sign), $S(fg,D,h)$ converges to $\iint_D f(x)g(y)\,\dd A$.
Taking the limit in the factorization gives
\[
    \boxed{\iint_D f(x)g(y)\,\dd A
    =\left(\int_a^b f(x)\,\dd x\right)
     \left(\int_c^d g(y)\,\dd y\right)}.
\]
\end{solution}

\end{document}
